\begin{frame}{체크리스트 적용: 5개 항목을 하나씩}
  \small
  \begin{tabular}{@{}p{2.9cm}p{11.7cm}@{}}
    \toprule
    항목 & 적용 \\
    \midrule
    문제 정의 & \textbf{명확} — 이진 분류, 행=환자, 라벨=악성/양성 \\
    방법론 적절성 & \textbf{여기가 핵심}. breast\_cancer는
      \textit{정형 표 데이터}. 발표팀의 답변(``복잡한 상호작용 자동
      학습'')은 16.1 FAQ가 명시적으로 허용한 근거 중 하나 —
      \textbf{형식상 통과}하나, FAQ가 요구한 ``\textbf{Ch07 GBDT와
      비교했을 때}''를 \textit{정량적으로} 충족했는가? ``GBDT 0.94
      vs 신경망 0.95''는 정확도로 \textbf{1\%p} 차이이고, 선택 편향
      계산($m$이 크면 부풀림 $\approx 0.05$)에 비하면
      \textbf{잡음보다 작은 차이} — ``신경망이 GBDT보다 좋다''는
      결론은 이 발표만으로 \textbf{성립하지 않음} \\
    수식적 정당화 & ``64개 히든 뉴런''이라는 \textit{숫자만} 있고,
      Ch09.1 파라미터 수 공식 ($30\times64 + 64 + 64\times2 + 2 =
      2114$개) 또는 Ch10.1 CNN 설계 공식으로 \textit{왜} 64인가를
      설명한 흔적이 없음 — 형식상 ``구조를 정의했다''는 것이지
      ``\textbf{정당화}한 것''이 아님 \\
    결과의 정직성 & ``잘 안 된 부분''이 발표에 \textbf{없음} —
      8.1 보고서 구조 §6이 \textit{필수}로 요구하는 절 \\
    ``왜 딥러닝인가'' & 답변은 있었으나(위 방법론 항목)
      \textbf{정량적 근거(GBDT와 val에서의 비교 \textit{표})}가
      없음 \\
    \bottomrule
  \end{tabular}
\end{frame}
